\section{\tkzClass{list\_point} 类}
\label{sec:class_list_point}

变量 \tkzVar{list\_point}{LP} 持有一个表，用于存放 \tkzClass{list\_point} 对象。
它的使用是可选的，用户可以自由选择任何变量名。

\medskip

不过，为了在整个文档中保持一致和可读，
推荐使用预定义变量 \tkzVar{list\_point}{LP}。
函数 \tkzFct{tkz-elements}{init\_elements()} 会自动重新初始化该表。

\medskip

\tkzClass{list\_point} 类表示一个有序的点列表。
它主要用作一个轻量容器，用于存放和处理由几何作图、交点或迭代算法产生的多个点。

与直线或圆等几何对象不同，\tkzClass{list\_point} 对象并不定义内在的几何。
它的目的是便于在 \tkzname{Lua} 与 \tkzname{TikZ} 之间进行分组、迭代和数据传递。

\medskip

一个 \tkzClass{list\_point} 对象的行为类似于一个以整数为索引的 Lua 表，
同时提供一组适合几何工作流的便捷方法。

\medskip

\subsection{属性}
\label{subsec:list_point_attributes}

一个 \tkzClass{list\_point} 对象只有一个属性 \tkzAttr{list\_point}{type}，其值为 \code{"list\_point"}。点的个数由 \tkzMeth{list\_point}{len()}（或 \code{\#LP.x}）给出，第 i 个点由 \tkzMeth{list\_point}{get(i)}（或 \code{LP.x[i]}）给出。

\subsection{创建对象}
\label{subsec:list_point_create}

\subsubsection{方法 \tkzMeth{list\_point}{new(...)}}
\label{ssub:meth_list_point_new}

创建一个新的 \tkzClass{list\_point} 对象。

在用户代码中，推荐使用可调用形式 \code{list\_point(...)} 作为构造器。
在内部，这会调用方法 \tkzMeth{list\_point}{new(...)}。

\medskip

一个 \tkzClass{list\_point} 对象既可以创建为局部 Lua 变量，
也可以作为表 \tkzVar{list\_point}{LP}（专门用于存放点列表）的一个条目。

\begin{tkzexample}[code only]
local P = list_point()
LP.Q = list_point(z.A, z.B, z.C)
\end{tkzexample}

两种形式都合法，并产生相互独立的 \tkzClass{list\_point} 对象。
参数必须是点；否则会引发错误。

\medskip

对于临时计算或中间结果，使用局部变量很方便。
不过，对于打算复用、传递给 \tkzname{TikZ} 或在局部作用域之外访问的对象，建议把它们存放在 \tkzVar{list\_point}{LP} 表中。

\medskip

函数 \tkzFct{tkz-elements}{init\_elements()} 会自动重新初始化
\tkzVar{list\_point}{LP} 表。


% --------------------------------------------------------------------
\subsection{方法或基本访问器}
\label{subsec:list_point_accessors}

\vspace{1em}
\bgroup
\small
\captionof{table}{\texttt{list\_point} 的方法。}\label{list_point:methods}
\begin{tabular}{ll}
\toprule
\texttt{方法}  & \texttt{参考} \\
\midrule
\texttt{构造器} &    \\
\midrule
\tkzMeth{list\_point}{new(...)}   & [\ref{ssub:meth_list_point_new}] \\
\midrule
\texttt{返回整数的方法}&\\
\midrule
\tkzMeth{list\_point}{len()}      & [\ref{ssub:meth_list_point_len}] \\
\midrule
\texttt{返回点的方法}&\\
\midrule\tkzMeth{list\_point}{get(i)}     & [\ref{ssub:meth_list_point_get}] \\
\tkzMeth{list\_point}{barycenter()}& [\ref{ssub:meth_list_point_barycenter}]\\
\midrule
\texttt{返回 List\_Point 的方法}&\\
\midrule
\tkzMeth{list\_point}{add(p)}     & [\ref{ssub:meth_list_point_add}] \\
\tkzMeth{list\_point}{extend(pl)} & [\ref{ssub:meth_list_point_extend}] \\
\tkzMeth{list\_point}{map(f)}     & [\ref{ssub:meth_list_point_map}] \\
\midrule
\tkzMeth{list\_point}{unpack()}   & [\ref{ssub:meth_list_point_unpack}] \\


\tkzMeth{list\_point}{clear()}    & [\ref{ssub:meth_list_point_clear}] \\
\tkzMeth{list\_point}{foreach(f)} & [\ref{ssub:meth_list_point_foreach}] \\

\tkzMeth{list\_point}{bbox()}     & [\ref{ssub:meth_list_point_bbox}] \\

\tkzMeth{list\_point}{to\_tikz\_path()}& [\ref{ssub:meth_list_point_to_path}] \\
\tkzMeth{list\_point}{as\_path()}& [\ref{ssub:meth_list_point_as_path}] \\
\bottomrule
\end{tabular}
\egroup

\subsubsection{方法 \tkzMeth{list\_point}{len()}}
\label{ssub:meth_list_point_len}

返回列表中存放的点的个数。

\subsubsection{方法 \tkzMeth{list\_point}{get(i)}}
\label{ssub:meth_list_point_get}

\code{get(i)} 返回索引为 \code{i} 的点；\code{get()} 返回所有点。

\subsubsection{方法 \tkzMeth{list\_point}{unpack()}}
\label{ssub:meth_list_point_unpack}

用 \code{table.unpack} 把所有点作为独立值返回。当需要与期望以逐个点作为参数的函数对接时，该方法很有用。

\begin{tkzexample}[latex=.5\textwidth]
\directlua{
init_elements()
LP.pts = list_point()
LP.pts:add(point(0,0))
LP.pts:add(point(1,0))
LP.pts:add(point(1,1))
z.A, z.B, z.C = LP.pts:unpack()
}
\begin{center}
\begin{tikzpicture}
\tkzGetNodes
\tkzDrawPolygon(A,B,C)
\end{tikzpicture}
\end{center}
\end{tkzexample}

% --------------------------------------------------------------------
\subsection{修改方法}
\label{subsec:list_point_mutators}

\subsubsection{方法 \tkzMeth{list\_point}{add(p)}}
\label{ssub:meth_list_point_add}

把点 \code{p} 追加到列表末尾，并返回对象本身。若 \code{p} 不是点则引发错误。

\subsubsection{方法 \tkzMeth{list\_point}{extend(pl)}}
\label{ssub:meth_list_point_extend}

把另一个 \tkzClass{list\_point} 对象 \code{pl} 中的所有点追加进来。若某个元素不是点则引发错误。

\subsubsection{方法 \tkzMeth{list\_point}{clear()}}
\label{ssub:meth_list_point_clear}

从列表中移除所有点。


这些方法就地修改对象，并支持方法链式调用。

% --------------------------------------------------------------------
\subsection{迭代辅助}
\label{subsec:list_point_iteration}

\subsubsection{方法 \tkzMeth{list\_point}{foreach(f)}}
\label{ssub:meth_list_point_foreach}

对点列表中的每个点 \code{p} 及其索引 \code{i} 应用函数 \code{f(p, i)}。

\subsubsection{方法 \tkzMeth{list\_point}{map(f)}}
\label{ssub:meth_list_point_map}

对每个点应用函数 \code{f(p, i)}，并返回一个新的 \tkzClass{list\_point} 对象，其中包含结果。该函数必须返回一个点。


这些方法为几何变换提供了一种函数式编程风格。

% --------------------------------------------------------------------
\subsection{几何工具}
\label{subsec:list_point_geometry}

\subsubsection{方法 \tkzMeth{list\_point}{barycenter()}}
\label{ssub:meth_list_point_barycenter}

返回列表中所有点的重心。
若列表为空，则该方法返回 \code{nil}。

\subsubsection{方法 \tkzMeth{list\_point}{bbox()}}
\label{ssub:meth_list_point_bbox}

以四个数返回点集的包围盒：
\[
  (x_{\min}, y_{\min}, x_{\max}, y_{\max})
\]


% --------------------------------------------------------------------
\subsection{TikZ 输出}
\label{subsec:list_point_output}

\subsubsection{方法 \tkzMeth{list\_point}{to\_tikz\_path()}}
\label{ssub:meth_list_point_to_path}

把点列表转换成形如下面的 TikZ 路径字符串：
\[
  (x_1,y_1) -- (x_2,y_2) -- \dots
\]
适合用 TikZ 命令绘图。

\subsubsection{方法 \tkzMeth{list\_point}{as\_path()}}
\label{ssub:meth_list_point_as_path}

返回一个包含该列表各点的 \tkzClass{path} 对象，可与路径宏配合使用（\code{\textbackslash tkzDrawCoordinates}、\code{\textbackslash tkzDrawPointsFromPath} 等）。见 [\ref{sub:macro_tkzDrawPointsFromPath}] 中的示例。


\medskip

这个类在 \tkzname{Lua} 计算与 \tkzname{TikZ} 绘图指令之间架起了一座简洁而有效的桥梁。


\subsection{示例}
\label{subsec:list_point_examples}

% --------------------------------------------------------------------
\subsubsection{临时点列表}
\label{ssub:list_point_example_local}

本例说明如何在一次计算中使用局部 \tkzClass{list\_point} 对象来存放临时结果。

\begin{tkzexample}[code only]
local P = list_point()
for i = 1, 10 do
  P:add(z.A + i * (z.B - z.A) / 10)
end
local G = P:barycenter()
\end{tkzexample}

列表 \code{P} 只在本地用于计算一组点的重心。

% --------------------------------------------------------------------
\subsubsection{把作图结果存放在 \tkzVar{list\_point}{LP} 中}
\label{ssub:list_point_example_lp}

本例说明如何存放一组交点供稍后使用。

\begin{tkzexample}[code only]
LP.I = list_point()
local X, Y = intersection(L.AB, C.circ)
LP.I:add(X)
LP.I:add(Y)
\end{tkzexample}

把点存放在 \tkzVar{list\_point}{LP} 中后，即可供绘图或进一步作图使用。

% --------------------------------------------------------------------
\subsubsection{迭代作图}
\label{ssub:list_point_example_iterative}

一个典型用例是逐步构造一族点。

\begin{tkzexample}[code only]
LP.P = list_point()
for i = 0, 24 do
  local t = i / 24
  LP.P:add(L.AB:point(t))
end
\end{tkzexample}

列表 \code{LP.P} 包含一条线段的离散采样。

% --------------------------------------------------------------------
\subsubsection{函数式变换}
\label{ssub:list_point_example_map}

本例使用 \tkzMeth{list\_point}{map} 方法来变换一组点。

\begin{tkzexample}[code only]
local P = list_point(z.A, z.B, z.C)
LP.Q = P:map(function(p)
  return z.O:rotation(tkz.tau / 3, p) rotation_(p, tkz.tau / 3, z.O)
end)
\end{tkzexample}

原始列表 \code{P} 保持不变，而 \code{LP.Q} 存放变换后的点。

% --------------------------------------------------------------------
\subsubsection{转换成 TikZ 路径}
\label{ssub:list_point_example_path}

本例说明 \tkzClass{list\_point} 与 TikZ 之间的交互。

\begin{tkzexample}[code only]
LP.P = list_point(z.A, z.B, z.C)
local path = LP.P:to_tikz_path()
\end{tkzexample}

\tkzMeth{list\_point}{to\_path} 返回的字符串可以直接传给 TikZ 绘图命令。


% --------------------------------------------------------------------
\subsubsection{从 \tkzClass{list\_point} 到 \tkzClass{path}，并用 \tkzNamePack{tkz-euclide} 绘图}
\label{ssub:list_point_example_points_path_draw}

本例展示一个典型工作流：在 \tkzname{Lua} 中计算一组点，
把它转换成一个 \tkzClass{path}，然后用 \tkzMeth{tkz-euclide}{tkzDrawPointsFromPath} 画出所得的点。


\begin{tkzexample}[latex = .5\textwidth]
\directlua{
  init_elements()
  z.A = point(0, 0)
  z.B = point(8, 6)
  L.AB = line(z.A, z.B)
  % list of sampled points
  LP.P = list_point()
  for i = 0, 12 do
    LP.P:add(L.AB:point(i/12))
  end
  % convert list_point -> path
  PA.sample = path()
  for i = 1, LP.P:len() do
    PA.sample:add_point(LP.P:get(i))
  end
}
\begin{center}
\begin{tikzpicture}
\tkzGetNodes
% draw the points stored in the path
\tkzDrawPointsFromPath(PA.sample)
\end{tikzpicture}
\end{center}
\end{tkzexample}

点存放在 \code{LP.P} 中以供稍后复用，而 \code{PA.sample} 包含对应的 TikZ 路径表示，用于绘图。
